\subsection{Forwarding and concurrency pilot: three measurement lessons}\label{sec:live-study}
A separate pilot holds initial messages fixed and varies whether supplementary tool records enter the next model request. \emph{Full} forwards the whole response; \emph{Projected} keeps only task-required fields. The synthetic records deliberately carry action information. Two fixed 8B models perform 3,120 local LangGraph/MCP executions over streaming TLS~1.3/HTTP/2. Each test condition has 120 tasks in twenty nonempty and four empty sessions. Supplement S1 reports every attack, phase, interval and outcome. The pilot is underpowered for its declared 10-point effect, so we report it for the evaluation lessons it supports and not as a causal result.

\paragraph{Lesson 1: single-step minimization does not contain the task}
First-feedback LR recognition falls from 43 to 23 correct for Qwen and from 40 to 22 for Llama. None of the eight primary local comparisons rejects after Holm correction (minimum adjusted $p=.0508$). Over the complete execution, the projected model view still reaches 44/120 and 33/120, against chance at 20/120. Later requests and responses carry task information that a first-step filter does not remove, which is Finding~1 measured on one deployment.

\paragraph{Lesson 2: shaping cost must be judged against concurrency}
In the concurrent conditions, each supplied window mixes the traffic of all tasks active at the same time. Recognition over both links is therefore already near chance with dedicated connections (30/120 and 20/120). Connection sharing (23 and 19) and joint model-link shaping (20 and 22) add no further detectable reduction. Joint shaping still costs 19.5--22.5$\times$ model-link bytes. It also leaves event localization unchanged: learned-detector F1 at IoU .5 is .117--.274 with shaping versus .077--.249 without. This echoes the classical observation that costly padding can buy little against coarse features~\cite{dyer2012peekaboo}.

\paragraph{Lesson 3: event metrics are sensitive to the matching rule}
Continuous detectors receive connection identities but no task windows. Their shared-condition precision is 4.3--18.6\%, with 8.21--44.79 false events per test session at IoU .5. A constant all-active baseline outranks the learned detectors at IoU .3 and falls below them at .5. Agent localization claims should therefore report several overlap thresholds and a constant-output reference.
